%

Wnioskiem z twierdzenia o dwusiecznej jest: % lang-pl
Una conseguenza del teorema della bisettrice è la seguente: % lang-it

\begin{theorem}
[Steinera-Lehmusa] % lang-pl
[di Steiner-Lehmus] % lang-it
\index{twierdzenie!Steinera-Lehmusa}% % lang-pl
\index{teorema!di Steiner-Lehmus}% % lang-it
\label{theorem_steiner_lehmus}%
    Jeżeli dwie dwusieczne trójkąta są równej długości, to trójkąt ten jest równoramienny. % lang-pl
	Se due bisettrici di un triangolo hanno la stessa lunghezza, allora il triangolo è isoscele. % lang-it
\end{theorem}

Po raz pierwszy wspomni o nim Christian Lehmus w~liście z~1840 roku do Charlesa Sturma, gdzie poprosi o~czysto geometryczny dowód. % lang-pl
Il teorema fu menzionato per la prima volta da Christian Lehmus in una lettera del 1840 a Charles Sturm, nella quale chiedeva una dimostrazione puramente geometrica. % lang-it
\index[persons]{Lehmus, Christian}%
\index[persons]{Sturm, Charles}%
Sturm przekazał prośbę do innych matematyków, jedną z~pierwszych osób, która uporała się z problemem, był Jakob Steiner. % lang-pl
Sturm trasmise la richiesta ad altri matematici; una delle prime persone a risolvere il problema fu Jakob Steiner. % lang-it
\index[persons]{Steiner, Jakob}%
Większość znanych uzasadnień to dowody nie wprost: jeśli trójkąt nie jest równoramienny, to ma dwusieczne różnej długości. % lang-pl
La maggior parte delle dimostrazioni note procede per assurdo: se un triangolo non è isoscele, allora le sue bisettrici hanno lunghezze diverse. % lang-it

Dowód można znaleźć u Hartshorne'a \cite[s. 11]{hartshorne2000}; Pompego \cite[s. 32, 33]{pompe_2022}; Bogdańskiej, Neugebauera \cite[s. 74]{neugebauer_2018}. % lang-pl
Zdzisław Pogoda w $\Delta_{11}^1$ napisze, jakie pomysłowe rozwiązanie znaleźli dwaj angielscy inżynierowie, G. Gilbert oraz D. McDonnell. % lang-pl
Una dimostrazione si trova in Hartshorne \cite[p. 11]{hartshorne2000}, Pompe \cite[pp. 32, 33]{pompe_2022} e Bogdańska, Neugebauer \cite[p. 74]{neugebauer_2018}. % lang-it
Zdzisław Pogoda, in $\Delta_{11}^1$, descrive l'ingegnosa soluzione trovata da due ingegneri inglesi, G. Gilbert e D. McDonnell. % lang-it
\index[persons]{McDonnell, D.}%
\index[persons]{Gilbert, G.}%
% TODO The American Mathematical Monthly (7,1963,str.79–80)
Coxeter \cite[s. 32]{coxeter_1967} podaje je w formie ćwiczenia (po tym jak wcześniej \cite[s. 26, 33]{coxeter_1967} poprosi o~dowód tego samego dla dwusiecznej zamienionej na środkową lub wysokość, co znacząco obniża poziom trudności). % lang-pl
Eves \cite[s. 19, 58]{eves1_1972} zachęca do dowodu z dopiskiem, że jest trudny. % lang-pl
Coxeter \cite[p. 32]{coxeter_1967} la propone come esercizio (dopo aver chiesto in precedenza \cite[pp. 26, 33]{coxeter_1967} una dimostrazione analoga con la bisettrice sostituita da una mediana o da un'altezza, il che riduce notevolmente la difficoltà). % lang-it
Eves \cite[pp. 19, 58]{eves1_1972} invita il lettore a dimostrarlo, osservando che non è facile. % lang-it

% https://www.algebra.com/algebra/homework/word/geometry/Medians-in-an-isosceles-triangle.lesson

%